% ==================== 2.4 小结表 ====================
表~\ref{tab:sym-summary} 汇总了本节六组对称元素演示的空间群、等效位置数与原子分布
特征。

\begin{table}[H]
  \centering
  \small
  \caption{六类对称元素演示的对比}
  \label{tab:sym-summary}
  \begin{tabular}{@{}lllccp{3.6cm}@{}}
    \toprule
    对称元素 & 国际符号 & 演示空间群（No.） & 一般位置数 & 平移分量 & 原子分布与辨识 \\
    \midrule
    反演中心 & $\bar1$ & $P\,\bar1$（2） & 2 & 无 & 关于反演中心等距反向的原子对 \\
    旋转轴 & $3$ & $P\,3$（143） & 3 & 无 & 绕轴 $120^\circ$ 分布，位于垂直轴的同一平面 \\
    旋转反演轴 & $\bar3$ & $P\,\bar3$（147） & 6 & 无 & 上下两层各 3 个，沿轴投影呈``似六重''分布 \\
    镜面 & $m$ & $P\,m\,\bar3$（200） & 24 & 无 & 三组互相垂直镜面，原子镜像成对 \\
    螺旋轴 & $3_1$ & $P\,3_1$（144） & 3 & $c/3$ & 沿轴螺旋上升，不位于同一平面 \\
    滑移面 & $c$ & $P\,1\,c\,1$（7） & 2 & $c/2$ & 反映后沿 $c$ 方向错开半个周期 \\
    \bottomrule
  \end{tabular}
\end{table}

把六组演示放在一起比较，可以看到一条清晰的分界：反演中心、旋转轴、旋转反演轴与
镜面属于宏观对称元素，产生的等效原子由纯旋转、反映或反演相联系，不含平移分量；
$3_1$ 螺旋轴与 $c$ 滑移面则分别带着 $c/3$ 与 $c/2$ 的平移，等效原子沿轴向逐个错开，
这种复合操作特征只有在超晶胞中才显示出来。六种演示基本覆盖了空间群国际符号中的
常见符号类型（旋转轴 $n$、旋转反演轴 $\bar n$、镜面 $m$、螺旋轴 $n_m$ 及各类滑移面
$a$、$b$、$c$、$n$、$d$），回头再看 $P\,\bar1$、$P\,3$、$P\,\bar3$、$P\,m\,\bar3$、
$P\,3_1$、$P\,1\,c\,1$ 这些符号，即可逐字读出它们的含义。对称元素是否含平移分量
（宏观还是微观），决定了单胞内等效原子的空间分布；而对称性同时决定衍射的系统消光
规律，这也是后续物相鉴定与空间群判定的实验基础。
